An AAC device speaks for a person, so a plausible but wrong sentence can misstate consent or refusal.
The gate lowers this risk without removing it: it checks selected concepts, polarity, and a limited set of intent cues, not meaning in full.
The evaluated interface speaks a sentence as soon as it is ready and does not show when a template replaced model output; the repaired build labels templates and does not speak a clarification or a literal preview unprompted, but consequential messages should still be confirmed with the user or a caregiver, and a preview-and-confirm step for generated sentences is a priority for deployments.
Direct \emph{Speak} and Quick Express remain available when generated text is unsuitable.

Voice cloning requires a reference clip and its exact transcript from the speaker or a guardian, with consent; clones are held in server memory only.
Optional features can send typed words to third parties (Google Input Tools, OpenAI Images) or store My Words and bookmarks in Firebase; generation, prediction, and speech do not need them.
The evaluation uses synthetic dialogues and scripted scenarios and contains no personal data; a study with AAC users would need accessible consent and the option to withdraw data.
Accessibility testing covered keyboard activation, emulated touch, and automated checks of names and live regions, not screen readers or switch scanning; the page has no live region for new text and renders nothing when the content delivery networks that serve its React, Babel, and Tailwind libraries are blocked, although generation and prediction run on the local server.
The evaluated build ran without the HMM-1 and HMM-2 checkpoints, so HMM-1 prediction and Un-Jumble are not evaluated.
